\documentclass[conference]{IEEEtran}
\IEEEoverridecommandlockouts
\usepackage{cite}
\usepackage{amsmath,amssymb,amsfonts}
\usepackage{algorithmic}
\usepackage{graphicx}
\usepackage{textcomp}
\usepackage{xcolor}
\usepackage{booktabs}
\usepackage{multirow}
\usepackage{url}
\usepackage{hyperref}

\def\BibTeX{{\rm B\kern-.05em{\sc i\kern-.025em b}\kern-.08em
    T\kern-.1667em\lower.7ex\hbox{E}\kern-.125emX}}

\begin{document}

\title{Baseline-Differential Memory Forensics: Diagnosing Capture-Environment Confounding in Automated Malware Detection Pipelines\\
% \thanks{}
}

% TODO(Joel): confirm final author block, affiliation, and whether Saltanov Kirill
% should appear as a co-author before submission.
\author{\IEEEauthorblockN{Joel C. Okore}
\IEEEauthorblockA{\textit{Faculty of Computer Science} \\
\textit{Innopolis University}\\
Innopolis, Tatarstan, Russia \\
j.okore@innopolis.university}

\and
\IEEEauthorblockN{Andrei Petrovski}
\IEEEauthorblockA{\textit{Faculty of Computer Science} \\
\textit{Innopolis University}\\
Innopolis, Tatarstan, Russia \\
a.petrovski@innopolis.ru}

\and
\IEEEauthorblockN{Igor V. Kotenko}
\IEEEauthorblockA{\textit{Security Problems Lab} \\
\textit{SPC RAS}\\
St. Petersburg, Russia \\
ivkote@comsec.spb.ru}
\textit{Institute of Information Security} \\
\textit{Innopolis University}\\
Innopolis, Tatarstan, Russia \\
i.kotenko@innopolis.ru
}

\maketitle

\begin{abstract}
Automated memory forensics pipelines increasingly pair real-time SIEM detection
with machine-learning classifiers trained on public memory-artefact datasets to
triage suspected malware infections without waiting for an analyst. We built and
validated such a pipeline, integrating Wazuh active response, WinPMEM memory
acquisition, Volatility3 analysis, and a gradient-boosted classifier, achieving
end-to-end detection and automated containment of a live ransomware infection in
a virtualised testbed. In validating the classifier, however, we found that its
99.8\% benchmark accuracy on the CIC-MalMem-2022 dataset does not reflect
malicious behavior: 22 of 22 count-based memory features are higher in the
\emph{benign} class than the malicious class, a hand-written zero-parameter rule
on a single feature matches the trained model's accuracy, and a propensity-score
analysis shows 99.08\% of the corpus is perfectly separable by capture protocol
alone, violating the positivity assumption required for any classifier trained on
it to generalise. We trace this to the dataset's own published methodology:
benign captures were taken from an active desktop and oversampled with SMOTE
\emph{before} the train/test split, while malicious captures came from a
2\,GB VirtualBox VM, so class label and capture environment are confounded by
construction. We propose a baseline-differential alternative that scores a
memory capture against the same host's own clean baseline rather than against
absolute counts, and validate it with a small paired pilot (six clean/infected
capture pairs from a controlled cloud VM, using MITRE ATT\&CK T1055 process-injection
techniques as a legal, reproducible surrogate for malware). The pilot surfaces a
mechanistic detectability boundary --- injection techniques that allocate private
executable memory are reliably flagged, while DLL-based and thread-hijacking
techniques are not --- and an independent case study in which a single
incompletely-cleaned injection contaminated every subsequent "clean" baseline for
the remainder of a capture session, itself direct evidence for the
capture-environment sensitivity this paper describes.
\end{abstract}

\begin{IEEEkeywords}
memory forensics, malware detection, dataset bias, machine learning, incident response, SOAR, Volatility3
\end{IEEEkeywords}

\section{Introduction}

Memory forensics is the analysis of a running system's volatile memory to recover
evidence --- processes, injected code, network state, hidden services --- that
disk-based analysis misses entirely, particularly for fileless malware and
in-memory rootkits. In an incident response setting, the value of a memory
capture decays quickly: the longer a compromised host runs after infection, the
more likely transient evidence is overwritten. This motivates automating the
acquisition step itself, triggering a memory dump the moment a SIEM raises a
credible alert rather than waiting for an analyst to request one.

We built such a pipeline. Wazuh, an open-source SIEM, monitors a Windows
endpoint for indicators such as encoded PowerShell execution or known malicious
cmdlets; a matching alert triggers a Wazuh active-response script that
remotely invokes WinPMEM over WinRM to acquire a memory image; the image is
analysed with Volatility3; a small set of extracted features is classified by a
pretrained model; and a positive classification triggers automated containment
(network isolation of the host) and alerting (a Slack notification with the
confidence score). We validated the complete pipeline end to end against a live
WannaCry ransomware infection in an isolated two-VM testbed, and it correctly
detected, captured, classified, and contained the infection in under three
minutes.

The classifier at the centre of that pipeline was trained on
CIC-MalMem-2022~\cite{carrier2022}, the standard public benchmark for
memory-artefact malware classification, and reproduced its reported accuracy of
approximately 99.8\%. In the course of hardening the pipeline for wider
deployment, we set out to understand \emph{why} two service-count features were
sufficient to reach that accuracy. What we found is the primary contribution of
this paper: the dataset's headline accuracy is an artefact of how it was
collected, not a measure of malware-detection capability, and the same failure
mode likely affects the substantial body of published work that reports similar
numbers on it.

This paper makes four contributions:

\begin{enumerate}
\item An automated, end-to-end memory-forensics response pipeline integrating a
SIEM, memory acquisition, Volatility3 analysis, and SOAR-style containment,
validated against a live ransomware infection (Section~\ref{sec:architecture}).
\item A quantitative diagnosis of capture-environment confounding in
CIC-MalMem-2022, including a positivity-violation analysis showing 99.08\% of
the corpus is separable by capture protocol alone, and a provenance trace to the
dataset's own published collection methodology (Section~\ref{sec:diagnosis}).
\item A baseline-differential feature representation that scores a capture
against the same host's own clean state rather than against absolute counts,
removing the confound by construction and requiring no malicious samples to
train (Section~\ref{sec:method}).
\item A paired pilot evaluation using MITRE ATT\&CK T1055 process-injection
techniques, which surfaces a mechanistic detectability boundary in the
\texttt{malfind} feature family and an independent, unplanned case study in
cross-capture contamination (Section~\ref{sec:evaluation}).
\end{enumerate}

\section{Related Work}

\textbf{Memory-artefact malware classification.} Carrier et al.~\cite{carrier2022}
introduced CIC-MalMem-2022 alongside a stacked-ensemble classifier reaching
99--99.02\% accuracy, extending VolMemLyzer~\cite{lashkari2019} feature
extraction with 26 obfuscation-focused features. The dataset has since been used
in numerous follow-on studies reporting comparably high accuracy with a range of
classical and deep-learning models, including work applying XGBoost and PCA-based
dimensionality reduction to the same corpus~\cite{abid2024}.

\textbf{SIEM-driven active response.} Wazuh's active-response mechanism, and
comparable capabilities in other open-source SIEM platforms, allow
detection rules to trigger arbitrary remote scripts, which prior work has used
for automated containment (e.g. host isolation, credential revocation) but --
to our knowledge -- not for triggering the memory-acquisition step of the
forensic pipeline itself, which is the integration this paper's system
performs.

\textbf{Dataset bias and evaluation validity in security ML.} The general
problem of classifiers learning spurious, environment-correlated shortcuts
rather than the intended semantic signal is well documented across security
machine learning, where non-representative or improperly split datasets are a
recurring cause of published results that do not replicate in deployment. Our
contribution is a concrete, quantified instance of this failure mode in a
specific and widely used memory-forensics benchmark, together with a remedy
that does not require a cleaner public dataset to exist.

\section{System Architecture}
\label{sec:architecture}

The pipeline spans two roles. A Windows 10 agent runs WinPMEM and is enrolled
with a Wazuh agent monitoring PowerShell operational logs and two watched
paths. An Ubuntu Wazuh manager runs the SIEM, custom detection rules, active
response scripts, Volatility3, and the classifier.

\subsection{Detection and Acquisition Trigger}
Custom Wazuh rules match indicators associated with MITRE ATT\&CK T1059.001
(PowerShell) and T1562.001 (defense evasion), including base64-encoded command
execution and known offensive cmdlets (e.g. \texttt{Invoke-Mimikatz}). A match
at severity level 10 or above fires an active-response command that invokes
\texttt{netexec} over WinRM to trigger \texttt{winpmem.exe} on the agent,
writing a raw memory image to a share visible to the manager.

\subsection{Feature Extraction and Classification}
File-integrity monitoring on the memory-dump path triggers a second active
response once the dump completes, running a Volatility3 plugin
(\texttt{windows.svcscan.SvcScan}) and extracting features into a CSV. A third
active response, triggered on modification of the features file, invokes the
classifier and acts on its output.

\subsection{Automated Response}
A positive classification triggers two actions: network isolation of the
infected host via a remote firewall command over WinRM, and a Slack alert
carrying the classification confidence and affected host. We validated this
path against WannaCry ransomware executed on the Windows agent via a
PowerShell download-and-run sequence; the pipeline detected, captured,
classified and isolated the host in under three minutes end to end.

% TODO(Joel): insert 1-2 figures here, e.g. Images/wazuh_alert_final.png and
% Images/response_and_pred_trig.png, with \includegraphics + \caption.

\section{Diagnosis: Capture-Environment Confounding}
\label{sec:diagnosis}

The classifier deployed in Section~\ref{sec:architecture} was trained on
CIC-MalMem-2022: 58{,}596 records, balanced 29{,}298 benign / 29{,}298
malicious, with 55 memory-artefact features per record extracted from
Windows memory images. Two features --
\texttt{svcscan.nservices} and \texttt{svcscan.process\_services} -- were
selected by feature importance and alone reach 99.84\% accuracy under a
random 70/30 split, matching the 55-feature model to within $10^{-4}$.

\subsection{A Zero-Parameter Baseline Matches the Trained Model}
A hand-written rule requiring no training --- classify as malicious if
$\text{nservices} \le 391$ or $\text{nservices} > 1000$ --- reaches 99.63\%
accuracy on the same corpus, statistically indistinguishable from the trained
model (Table~\ref{tab:ablation}).

\begin{table}[htbp]
\caption{Feature-Group Ablation, Random 70/30 Split}
\label{tab:ablation}
\centering
\begin{tabular}{lcc}
\toprule
Feature set & \# features & Accuracy \\
\midrule
All 55 features & 52 & 0.9999 \\
Host-configuration counters only & 6 & 0.9998 \\
Injection/hiding features only & 24 & 0.9998 \\
Deployed model (2 \texttt{svcscan} features) & 2 & 0.9984 \\
Zero-parameter threshold rule & 0 & 0.9963 \\
\bottomrule
\end{tabular}
\end{table}

\subsection{Effect Directions Are Physically Inverted}
Across all 22 count-based features in the dataset, the benign class has a
\emph{higher} median value than the malicious class in every single case, with
no exceptions. This includes the two canonical injection and process-hiding
indicators: \texttt{malfind.ninjections} (benign median 4, malicious median 3)
and \texttt{ldrmodules.not\_in\_load}, a hidden-module count (benign median 74,
malicious median 46). Under any behavioral hypothesis these should point the
opposite way.

\subsection{No Common Support}
We fit a five-fold cross-validated logistic-regression propensity score
$e(x) = P(\text{malicious} \mid x)$ over all 55 features. Only 0.92\% of the
corpus falls in the overlap band $0.1 < e(x) < 0.9$; 99.08\% is perfectly
separated by the classifier with no ambiguity whatsoever
(Fig.~\ref{fig:overlap}). Load-matching benign and malicious captures on a
system-load proxy (\texttt{handles.nhandles}) leaves 148 of 58{,}596 rows in
common support: the 5th percentile of the benign distribution already exceeds
the 95th percentile of the malicious distribution. This is a positivity
violation in the formal sense used in causal inference: when treatment (here,
class label) and confound (capture environment) are this tightly coupled, no
reweighting, matching, or choice of classifier can recover a causal, i.e.
behavioral, estimate from the data as collected.

\begin{figure*}[htbp]
\centering
\includegraphics[width=\textwidth]{fig_overlap.png}
\caption{(a) Propensity score distribution showing near-complete separation.
(b) Disjoint system-load distributions between classes. (c) Least-overlapping
individual features.}
\label{fig:overlap}
\end{figure*}

\subsection{Provenance}
The dataset's own methodology paper explains the mechanism directly. Benign
captures were drawn from "normal user behavior \ldots using different
applications in the machine" and then \emph{"performed oversampling using SMOTE
algorithm to make the dataset balanced"}~\cite{carrier2022}, applied before any
train/test split existed, meaning synthetic benign rows and the real rows used
to generate them can and do land on opposite sides of a random split. We
confirmed this independently: \texttt{malfind.uniqueInjections} takes 209
distinct values in the malicious class, all small-integer ratios consistent
with genuine counting (1.0, 1.1, 1.125, \ldots), while the benign class takes
5{,}337 distinct values densely packed as irrational-looking decimals (e.g.
1.000433596), consistent with SMOTE interpolation rather than measurement.
Separately, malicious samples were executed "in a VM with 2 Gigabytes of
memory"~\cite{carrier2022} while benign captures came from an active desktop;
a 2\,GB VM structurally cannot host as many processes, handles, and services as
a working desktop, independent of any malware behavior, which is sufficient
on its own to explain the inverted effect directions in Section~IV-B.

To our knowledge this diagnosis has not previously been reported; existing
work on CIC-MalMem-2022 addresses feature selection and adversarial poisoning,
not the collection protocol itself.

\section{Baseline-Differential Method}
\label{sec:method}

Given that absolute memory counters are dominated by host configuration and
workload rather than infection state, we propose scoring a capture against the
distribution of the \emph{same host's own clean captures} rather than against
a population-level threshold.

For a feature vector $x$ extracted from a candidate capture and a set of clean
reference captures from the same host, we compute a robust deviation
\begin{equation}
z_i = \frac{x_i - \text{median}_i}{1.4826 \cdot \text{MAD}_i + \epsilon}
\end{equation}
where $\text{MAD}_i$ is the median absolute deviation of feature $i$ over the
clean reference set, floored at $\max(0.01\,|\text{median}_i|,\,1)$ so that a
feature with near-zero baseline variance does not produce an unbounded score
from a single unit of change. A companion scale-free ratio
$r_i = (x_i+1)/(\text{median}_i+1)$ is retained for interpretability. The
model requires \emph{only clean captures to fit}: no malicious samples are
needed for training, which sidesteps the labelled-malware scarcity that drives
datasets like CIC-MalMem-2022 toward synthetic augmentation in the first
place.

\subsection{Paired Collection Protocol}
To validate the method without inheriting the confound it corrects, we designed
a within-subject collection protocol: every infected capture is paired with a
clean capture from the \emph{same} virtual machine, same snapshot lineage,
within tens of minutes, with the intervening workload varied across a small
rotation (idle, browser activity, office-application activity, mixed, and a
compute-bound "busy" state) so the clean baseline itself carries realistic
variance rather than the near-zero dispersion observed in
Section~\ref{sec:diagnosis}.

\section{Pilot Evaluation}
\label{sec:evaluation}

Full malware execution requires an isolated lab we did not have available for
this iteration of the work. As a legal, reproducible, and mechanistically
transparent surrogate, we used Atomic Red Team~\cite{atomicredteam}
implementations of five MITRE ATT\&CK T1055 (process injection) sub-techniques
on a single Windows Server 2022 cloud VM: T1055.001 (DLL injection via
\texttt{mavinject.exe}), T1055.002 (portable executable injection),
T1055.003 (thread execution hijacking), T1055.004 (APC queue injection), and
T1055.012 (process hollowing).

\subsection{Data Collection Integrity}
Ten paired trials were executed sequentially. Inspecting the resulting
timestamps and feature trajectories revealed that from the seventh trial
onward, a memory-residency artefact from an earlier T1055.001 injection
persisted across the process-level cleanup performed between trials --
Atomic Red Team's cleanup routines remove dropped files and registry
modifications but cannot retroactively un-inject a DLL already resident in a
still-running process -- inflating \texttt{malfind.commitCharge} from a stable
baseline of 259--265 to a new stable baseline of 773--778 for every subsequent
capture, clean and infected alike, for the remainder of the session
(Table~\ref{tab:contamination}). We treat the first six trials, prior to this
event, as the valid paired dataset ($n=6$) for the quantitative results below,
and the remaining four trials as an independent case study.

\begin{table}[htbp]
\caption{Cross-Trial Contamination Event}
\label{tab:contamination}
\centering
\begin{tabular}{lcc}
\toprule
Trial range & \texttt{malfind.commitCharge} & Status \\
\midrule
1--6 (all captures) & 259--265 & stable \\
7--10 (all captures) & 773--778 & stable, elevated \\
\bottomrule
\end{tabular}
\end{table}

This event is itself evidence supporting this paper's central claim: a single
incompletely-reversible behavioral artefact was sufficient to shift a memory
counter by nearly $3\times$ and hold it there indefinitely, on the very same
host, with no change in the ground-truth infection state at capture time. A
capture-environment confound of this kind, introduced across sequential
samples rather than between two different machines, is arguably harder to
detect than the one diagnosed in Section~\ref{sec:diagnosis} because within-host
provenance can create a false sense of safety; we recommend a full reboot,
not merely process termination, between captures in any future collection
protocol using this method.

\subsection{Paired Deltas}
Table~\ref{tab:deltas} reports the infected-minus-clean delta for each valid
pair on features associated with code injection and module hiding.
\texttt{ldrmodules.not\_in\_load} moves in the expected (positive) direction in
five of six pairs; the \texttt{malfind} feature family moves positively or
remains flat in all six pairs and never reverses.

\begin{table}[htbp]
\caption{Infected$-$Clean Deltas, Valid Pairs ($n=6$)}
\label{tab:deltas}
\centering
\resizebox{\columnwidth}{!}{%
\begin{tabular}{lccccc}
\toprule
Technique & \texttt{malfind.} & \texttt{malfind.} & \texttt{ldr.not\_} & \texttt{ldr.not\_} \\
 & \texttt{ninjections} & \texttt{commitCharge} & \texttt{in\_load} & \texttt{in\_mem} \\
\midrule
T1055.001 (idle)   & +1 & +2 & +1  & +1  \\
T1055.002 (browser)& 0  & 0  & --2 & --2 \\
T1055.003 (office) & 0  & 0  & +1  & +1  \\
T1055.004 (mixed)  & +3 & +3 & +15 & +15 \\
T1055.012 (busy)   & +1 & +7 & +5  & +5  \\
T1055.001 (idle, repeat) & 0 & 0 & +21 & +21 \\
\bottomrule
\end{tabular}}
\end{table}

\subsection{A Mechanistic Detectability Boundary}
The pattern in Table~\ref{tab:deltas} is not random: it separates cleanly
along the injection mechanism each technique uses. \texttt{malfind} is
designed to flag memory regions that are private, anonymous, and executable --
the signature of raw shellcode written directly into a process. T1055.004 and
T1055.012, whose Atomic Red Team implementations call
\texttt{VirtualAllocEx}, \texttt{WriteProcessMemory}, and
\texttt{VirtualProtectEx} to place and execute shellcode, are the two
techniques with a strong and consistent signal. T1055.001 (DLL injection),
T1055.002 (portable-executable injection into a file-backed section), and
T1055.003 (thread hijacking, which reuses an existing thread rather than
allocating new memory) leave \texttt{malfind}-family features unchanged in
most or all trials, consistent with this feature family's well-known blind
spot for injection techniques that do not rely on private anonymous
executable memory.

We repeated the baseline-differential scoring of Section~\ref{sec:method}
restricted to the eleven \texttt{malfind} and \texttt{ldrmodules} features.
Unlike a score computed over all 44 extracted features -- which, at $n=6$,
produced a false positive on a clean capture driven by unstable
median-absolute-deviation estimates on noisy host-configuration features such
as \texttt{svcscan.nactive} -- the behaviorally-restricted score produced no
false positives on any clean capture and correctly ranked the infected capture
above its paired clean capture for T1055.004 and T1055.012, the two
techniques Section~VI-C identifies as detectable by this feature family.

\section{Limitations}

The pilot is small ($n=6$ valid pairs after excluding the contamination
episode) and drawn from a single host image (Windows Server 2022), which was
itself a deployment constraint rather than a design choice: the standard
Windows 10/11 client images were unavailable on the student cloud
subscriptions used for this work without additional licensing attestation.
Atomic Red Team techniques exercise a single MITRE ATT\&CK sub-technique in
isolation and are not equivalent to a full malware kill chain involving
persistence, command-and-control, and exfiltration; the WannaCry validation in
Section~\ref{sec:architecture} demonstrates the pipeline handles real malware,
but that validation predates the baseline-differential feature set and should
be repeated once lab access to real malware samples is available. The
all-feature composite anomaly score requires a larger clean reference set
before its per-feature deviation estimates are stable; we report the
behaviorally-restricted variant for this reason and regard the full-feature
score as future work rather than a validated result.

\section{Conclusion}

We diagnosed a severe, previously unreported capture-environment confound in
CIC-MalMem-2022, tracing it to the dataset's own collection methodology and
showing that 99.08\% of the corpus is separable by capture protocol alone,
independent of infection state. We proposed a baseline-differential
alternative that scores a capture against the same host's own clean state, and
validated it in a small controlled pilot that surfaced both a genuine,
mechanistically explicable detection signal and an independent case study in
cross-capture contamination that reinforces the paper's central finding. Future
work includes a larger, reboot-isolated paired corpus across multiple host
images, restoring real malware validation once lab access is available, and
integrating the baseline-differential scorer into the deployed Wazuh pipeline
described in Section~\ref{sec:architecture} in place of the confounded
classifier it currently uses.

\begin{thebibliography}{00}

\bibitem{carrier2022} T. Carrier, P. Victor, A. Tekeoglu, and A. H. Lashkari,
``Detecting Obfuscated Malware Using Memory Feature Engineering,'' in
\textit{Proc. 8th Int. Conf. Information Systems Security and Privacy
(ICISSP)}, 2022, pp. 177--188.

\bibitem{lashkari2019} A. H. Lashkari, B. Li, T. L. Carrier, and P. Kaur,
``VolMemLyzer: Volatile Memory Analyzer for Malware Classification using Feature
Engineering,'' in \textit{Proc. Reconciling Data Analytics, Automation,
Privacy, and Security: A Big Data Challenge (RDAAPS)}, 2019.

\bibitem{abid2024} M. N. Abid, M. Beggas, and A. Laouid, ``Machine/Deep
Learning for Obfuscated Malware Detection,'' in \textit{Proc. 7th Int. Conf.
Future Networks and Distributed Systems (ICFNDS)}, 2024, pp. 390--396,
doi: 10.1145/3644713.3644768.

\bibitem{atomicredteam} Red Canary, ``Atomic Red Team,'' [Online]. Available:
\url{https://github.com/redcanaryco/atomic-red-team}

% TODO(Joel): verify every entry above against the publisher record before
% submission (page numbers, venue names, DOIs). The Carrier and Abid entries
% were extracted directly from source documents in this project; Lashkari and
% Atomic Red Team should be double-checked. Add: Volatility3 documentation,
% Wazuh active-response documentation, MITRE ATT&CK T1055 reference.

\end{thebibliography}

\end{document}
